% =====================================================================
% Section 7 --- Conclusion (merged draft; supersedes (7)discussion.tex
% and the stub (8)conclusion.tex)
% ---------------------------------------------------------------------
% Job: return to the thesis; position against BF as the frictionless pole;
%      carry the closure lineage (Sen, Taylor, Rose, FIDELIO); state the
%      limitations; close with practical guidance. NO bridge-discovery
%      claims -- the bridge is measurement. Kuhnian frame kept minimal.
% Mindmap: N03, N02D.
% Status: first full draft (2026-10-22). Absorbs the two paragraphs the
%      dropped Discussion would have carried (positioning, limitations).
% RETIRED (2026-10-24): "the Delta Pi/Y range from the financing wedge (Sec. 5.2)".
%      No financing-wedge section was written -- Sec. 5 carries only aggregate and
%      sectoral results -- and the number does not exist: the kernel prices at
%      marginal cost, so Delta Pi/Y is identically 0.00 in every plain financing
%      cell and is tabulated only once a pricing extension is switched on (Tables 7
%      and E). Do not re-open without a new run generation.
% =====================================================================
\section{Conclusion}\label{sec:conclusion}

Input--output models of a socio-ecological transformation usually derive multipliers, employment effects and fiscal bills from the Leontief inverse once final demand is given. For estimating a transformative demand shock, however, the inversion itself, which environmentally extended applications rightly inherit, matters less than the choice of closure made before it. This choice requires considering competing theories of what the short run is and which mechanisms underpin our monetary system and financing conditions. Moreover, it can move the aggregate result by tenths of a percentage point across most of the matrix, and yields a multiplier above one only in the interaction that pairs a pinned wage with external finance.

Our contribution is an illustrative comparison with a practical purpose. One core framework substantiates fifteen scenarios whose differences are entirely attributable to the assumptions that define them. Read this way, the matrix shows where the closure choice matters most: in the magnitude of the aggregate effect (the two extremes sit in the employment-adjustment and Leontief columns) and in whether sectoral relative prices move at all (the sectoral models).

This result is consistent with the existing literature. The accounting with intermediate inputs of \citet{Hulten.1978} and its modern descendants \citep{Baqaee.2021,Acemoglu.2012} describe the frictionless pole of our matrix --- the standard-CGE limit in which labor is perfectly mobile across sectors at one flexible wage, aggregate employment is given, and reallocation carries no cost. Our five closures represent  canonical answers of a long debate --- \citet{Sen1963} on the demarcation of neoclassical and neo-Keynesian closures, \citet{Taylor1990} on structuralist closures, \citet{Rose.1995} on their translation into computable form ---
carried into an accounting structure that transformation research already
uses.

Three limitations apply to the claims. First, labor is the only production factor in the model, so our results speak to the labor-market channel of the short run and say little about the capital- and investment-driven side of adjustment. Nor is the factor itself differentiated: a skill differential would allow the same question to be asked about the composition rather than the level of labor demand, which is an established extension that can enhance validity if specific skill shortages are expected for a given scenario. Second, the five welfare measures are dispersion-blind: none of them says anything about who within the household sector bears the adjustment. Third, the horizon is static: we compare short-run equilibria, not adjustment paths. An interesting vantage point for a dynamic reinterpretation of our efforts is given by \citet{McGregor.1996}, who shows how a capacity-constrained economy, where investment responds to demand, will converge into a model with demand-determined quantities.

For the applied modeler and transformation researcher, our paper provides a first, illustrative menu of closure choices and a corresponding open toolbox, \href{https://github.com/theplatters/OpenCGE}{OpenCGE}, to put them into practice. For applied work we hope to inspire more reflection about the interplay of foundational assumptions, closure choice for research on a socio-ecological information to better exploit the nuances and variations in possible model setups to articulate specific transformation scenarios and to more concisely illustrate under which scenarios should scenarios could succeed to truly improve both, sustainability as well as prosperity.
